% 本文件由 LaTeX 排版 Agent 根据有效 translation.json 自动生成。
% author=Augustine of Hippo; work_id=1407; title=论善的本性

\chapter{论善的本性}

% structure: p0001 source=h1 target=omitted reason=page-title-already-rendered-by-parent

写于404年之后。在\WorkTitle{订正录}中，它紧接在\WorkTitle{与摩尼教徒费利克斯论辩录}之后，后者写于404年底左右。这是最富论证性的反摩尼教论著之一，因而也是最深奥、最难的论著之一。这里所采取的论证思路，在较早的论著中已部分使用。论著中最有趣且对摩尼教徒最具杀伤力的部分，是取自摩尼的\WorkTitle{宝藏}和他的\WorkTitle{基本书信}的大段引文。——\Signature{A.H.N.}\par

% structure: p0003 source=h2 target=section reason=relative-heading-rank; base=h2

\section{第一章——天主是最高的、不变之善，一切其他善事（精神的与形体的）皆由祂而来}

至善，无高于其上者，即是天主；因此祂是不变之善，故而是真正永恒、真正不朽的。一切其他善事仅由祂而来，并非出于祂自身。因为出于祂自身的，即是祂自己。因此，既然唯有祂是不变的，那么祂所造的一切，因祂是从无中造出，故而是可变的。因为祂是如此全能，甚至能从无中，即从绝对不存在者中，造出或大或小的善事，或属天或属地，或精神或形体。但因祂也是公义的，祂并未使祂从无中所造之物与祂从自身所生者平等。因此，无论大小的善事，经由各种等级的受造物，只能来自天主；但因每个本性，就其是本性而言，皆是善的，故无任何本性可以存在，除非来自至高的真天主：因为一切事物，即使不是至善，却与至善相关；再者，因为一切善事，即使是最低等的、远离至善的，也只能从至善获得存在。因此，每一个精神，虽是可变的，以及每一个形体实体，都来自天主；而这一切被造之物，都是本性。因为每一个本性或是精神，或是形体。不变的精神是天主；可变的精神是被造的本性，但优于形体；而形体不是精神，除非当风——因其对我们不可见，而其力量却被感知为不小——在某种意义上被称为精神时。\par

% structure: p0005 source=h2 target=section reason=relative-heading-rank; base=h2

\section{第二章——此事如何足以纠正摩尼教徒}

但是，为了那些不能理解一切本性（即每一个精神和每一个形体）在本性上是善的，反而因精神的邪恶和形体的可朽而动摇，并因此企图引入另一种由天主所不造的邪恶精神和可朽形体的本性的人，我们决意如此引导他们理解我们所说能理解的事。因为他们承认，没有善事能存在除非来自至高的真天主，这既是真理，也足以纠正他们，只要他们愿意留心。\par

% structure: p0007 source=h2 target=section reason=relative-heading-rank; base=h2

\section{第三章——度量、形式与秩序：天主所造之物中的普遍善}

因为我们天主教基督徒敬拜天主，一切善事无论大小都从祂而来；一切度量无论大小都从祂而来；一切形式无论大小都从祂而来；一切秩序无论大小都从祂而来。因为一切事物，越是被更好地度量、形成和有序化，就确定无疑地更善；而越是被次要地度量、形成和有序化，就越不善。因此，这三者——度量、形式和秩序——（且不说无数其他归于这三者的事物）——这三者，即度量、形式和秩序，可以说是天主所造之物（无论是在精神中还是在形体中）的普遍善。因此，天主超越受造物的每一个度量、每一个形式、每一个秩序，祂并非以空间位置超越，而是以不可言喻的、独特的权能超越；一切度量、一切形式、一切秩序都来自祂。这三者，在哪里是大的，就是大善；在哪里是小的，就是小善；在哪里它们缺席，就没有善。同样，在哪里这三者是大的，就有大的本性；在哪里是小的，就有小的本性；在哪里它们缺席，就没有本性。因此，一切本性都是善的。\par

% structure: p0009 source=h2 target=section reason=relative-heading-rank; base=h2

\section{第四章——恶是对度量、形式或秩序的腐化}

因此，当追问恶从何而来时，必须先问恶是什么。恶不是别的，正是属于本性的度量、形式或秩序的腐化。因此，已被腐化的本性被称为恶；因为确实，当它未腐化时，它是善的；但即使腐化时，就它仍是本性而言，它是善的；就它被腐化而言，它是恶的。\par

% structure: p0011 source=h2 target=section reason=relative-heading-rank; base=h2

\section{第五章——有时，一个更优越等级的被腐化本性，甚至优于一个低等的未腐化本性}

但可能发生这样的情况：一个因本性度量与形式而被列为更优越的本性，即使腐化了，仍优于另一个因低下本性度量与形式而被列为低等的未腐化本性：如人按呈现于眼前的品质来评价，腐化的黄金无疑优于未腐化的白银，腐化的白银优于未腐化的铅；同样，在更有力的精神本性中，一个因邪恶意志而腐化的理性精神优于一个虽未腐化却无理性的精神；而任何精神，即使腐化，也优于任何形体，即使未腐化。因为一个本性，当它寓于身体时，为身体提供生命，这比那个被提供生命的本性更优。然而，纵然被造的生命精神可能已腐化，它仍能为身体提供生命，因此，即使腐化，它也优于未腐化的身体。\par

% structure: p0013 source=h2 target=section reason=relative-heading-rank; base=h2

\section{第六章——不能腐化的本性是至善；能腐化的本性也是某种善}

但如果腐化从可腐化之物中夺去一切度量、一切形式、一切秩序，就没有本性留存。因此，每一个不能腐化的本性是至善，正如天主那样。但每一个能腐化的本性本身也是某种善；因为腐化不能伤害它，除非夺取或减少那属于善的东西。\par

% structure: p0015 source=h2 target=section reason=relative-heading-rank; base=h2

\section{第七章——理性精神的腐化，一方面是自愿的，另一方面是惩罚性的}

但是，对于最卓越的受造物，即理性受造物，天主赐予了这一点：如果他们不愿意，就不会被腐坏；也就是说，如果他们能在主、他们的天主之下保持顺服，就会依附于祂不朽的美；但如果他们不愿保持顺服，既然他们甘愿在罪中腐坏，就将不情愿地在惩罚中腐坏，因为天主是如此美善，以至于任何离弃祂的人都不会有福，而在天主所造的万物中，理性的本性是如此伟大的美善，以至于除了天主之外，没有什么美善能使之真福。因此，罪人被预定受罚；这预定之所以是惩罚，是因它不符合他们的本性，但它是正义的，因为它符合他们的过错。\par

% structure: p0017 source=h2 target=section reason=relative-heading-rank; base=h2

\section{从低劣之物的腐坏与毁灭中生出宇宙之美}

至于其余由无中受造之物，它们无疑低于理性灵魂，既不能享福，也不能受苦。但由于事物本身按其形态与显相是善的，而且除非来自天主，不可能有较小或最小的善，它们如此被安排，使得较脆弱的让位于较坚固的，较软弱的让位于较强壮的，较无能的让位于较有力的；这样，尘世之物与天上之物和谐一致，居从于卓越者。然而，对于消逝与更替的事物，在其类中具有某种暂时的美，使得那些死去或不再成为原来状态的事物，并不降低或扰乱普遍受造的形态、显相与秩序；正如一篇措辞优美的演说，虽然其中的音节与声音仿佛在生与灭中疾驰而过，但它无疑是美丽的。\par

% structure: p0019 source=h2 target=section reason=relative-heading-rank; base=h2

\section{为有罪的本性设立的惩罚使之被正当地安排}

什么样的惩罚，惩罚何等重大，与每个过错相当，这属于天主的审判，而非人的审判；这惩罚，当它被归于皈依者而得以赦免时，天主方面表现出极大的良善；当它应得地被施加时，天主方面并无不义；因为本性在惩罚中正义地遭受痛苦，比在罪中不受惩罚地欢喜，更能被妥善安排；这本性，即便如此，仍有某种度量、形态与秩序，在任何极端情况下仍有一些善；如果这些善被完全拿走、彻底消灭，那么就没有善了，因为没有本性会存留。\par

% structure: p0021 source=h2 target=section reason=relative-heading-rank; base=h2

\section{可腐坏的本性，因由无中受造}

因此，一切可腐坏的本性之所以为本性，仅仅在于它们\Emph{来自}天主；如果它们是\Emph{出于}祂，就不会是可腐坏的，因为它们就会是祂自身所是。所以，不论它们有什么样的度量、什么样的形态、什么样的秩序，它们之所以如此，是因为造它们的是天主；但它们是可变的，因为它们是由无受造的。因为，把无与天主等同，是亵渎的鲁莽，如同我们想要使生于天主的，与祂由无中所造的一样。\par

% structure: p0023 source=h2 target=section reason=relative-heading-rank; base=h2

\section{天主不能受损害，也没有任何本性能在未得祂许可的情况下受损害}

因此，天主的本性既不能受损害，在天主之下的任何本性也不能不义地受损害：因为当某些人因不义地犯罪而加害时，不义的意志归于他们；但他们被许可加害的能力只来自天主，祂知道他们应当受什么苦，而他们自己却不知道，祂许可他们加害于谁。\par

% structure: p0025 source=h2 target=section reason=relative-heading-rank; base=h2

\section{一切善的事物唯独来自天主}

所有这一切如此清晰、如此确定，以至于那些引进另一种不由天主所造的本性的人，如果愿意留心，就不会充满如此巨大的亵渎，即把如此伟大的善置于至高的恶中，又把如此巨大的恶置于天主内。因为真理迫使他们承认这一点，即一切善的事物唯独来自天主，这足以纠正他们，如果愿意留心，如我上面所说。因此，并非大的善来自一位，小的善来自另一位；而是大大小小的善全都来自那至善的独一者，即天主。\par

% structure: p0027 source=h2 target=section reason=relative-heading-rank; base=h2

\section{每一个别的善事物，无论大小，都来自天主}

因此，让我们在心中呈上无论多大的善事物，这些善事物都宜归于天主作为其创造者；并将这些善事物除去之后，让我们看看是否还有什么本性存留。一切生命，无论大小；一切能力，无论大小；一切安全，无论大小；一切记忆，无论大小；一切美德，无论大小；一切理智，无论大小；一切安宁，无论大小；一切丰裕，无论大小；一切感觉，无论大小；一切光明，无论大小；一切甘饴，无论大小；一切度量，无论大小；一切美，无论大小；一切和平，无论大小；以及其他任何类似的事物，尤其是那些贯穿于一切事物中的，无论是精神性的还是形体性的，每一个度量、每一个形态、每一个秩序，无论大小，都来自主天主。凡愿意滥用所有这些善事物的人，将受天主的审判而受罚；但若这些事物完全都不存在，就没有任何本性存留。\par

% structure: p0029 source=h2 target=section reason=relative-heading-rank; base=h2

\section{较小的善与较大的善相比，被称为相反的名称}

但在所有这些事物中，凡是较小的，与较大的相比，就被称为相反的名称；例如在人的形象中，因为美更大，猿猴的美与之相比被称为畸形。愚昧的人被欺骗，仿佛前者是善，后者是恶；他们并不留意猿猴的身体中自身的形态、两侧肢体相等、各部协调、保护安全，以及其他不胜枚举的事物。\par

% structure: p0031 source=h2 target=section reason=relative-heading-rank; base=h2

\section{在猿猴的身体中，美之善是存在的，只是程度较低}

但为使我们所说的话能被理解，并能满足那些理解力过于迟钝的人，甚至使那些顽梗悖逆、抗拒最明显真理的人也被迫承认真理，可以问他们：腐朽能否损害猿猴的身体？如果能，使其变得更加丑陋，那么它所减损的不正是美的善吗？因此，只要身体的本性存续，就会保留一些东西。因此，如果善被耗尽，本性也就被耗尽，所以本性是善的。同样，我们说慢是快的对立面，但完全不动的甚至不能被称为慢。我们说低沉的声音与尖锐的声音相对，或刺耳的声音与悦耳的声音相对；但如果你完全除去任何声音，便有沉默，那里没有声音，而这种沉默，仅仅因为没有声音，通常被当作与声音相对立的东西。同样，光明与黑暗被称为仿佛两个相反的东西，但甚至黑暗也有某些光，当光完全不存在时，黑暗就是光的缺乏，正如沉默是声音的缺乏一样。\par

% structure: p0033 source=h2 target=section reason=relative-heading-rank; base=h2

\section{第十六章——事物中的缺乏由天主适当地安排}

然而，这些事物中的缺乏在万有的自然中如此被安排，以至于对明智思考的人来说，它们并非不适当地有其变迁。因为天主通过不照亮某些地方和时刻，也恰当地造成了黑暗，如同白昼一样。因为如果我们通过压制声音，在说话中恰当地插入沉默，那么祂，作为万有的完美创造者，岂不更恰当地造成事物中的缺乏吗？因此，在\WorkTitle{三少年歌}中，光明与黑暗一同赞美天主（\BibleRef{达 3:72}），即在善加思考的人心中引出赞美。\par

% structure: p0035 source=h2 target=section reason=relative-heading-rank; base=h2

\section{第十七章——本性，就其为本性而言，并非恶}

因此，没有任何本性，就其为本性而言，是恶的；但对每个本性而言，除了在善方面受到减损之外，并无恶。但如果因减损而被耗尽，以至于没有善留下，则无任何本性存留；不仅摩尼教徒所引入的那种本性（其中发现了如此多的善，以至他们的极度盲目是令人惊奇的），而且任何人所能引入的本性，都是如此。\par

% structure: p0037 source=h2 target=section reason=relative-heading-rank; base=h2

\section{第十八章——古人称为万物无形质料的\OriginalTerm{原质}{Hyle}并非恶}

因为就连那种古人称为\OriginalTerm{原质}{Hyle}的质料，也不可称为恶。我并非指摩尼以最无谓的虚妄（不知所说的是什么）所称的\OriginalTerm{原质}{Hyle}，即有形体的创造者；因此，人们正确地对他指出，他引入了另一个神。因为除了天主独自，无人能形成并创造有形体的存在；因为除非它们同时具有度量、形式和秩序，它们才能被创造，而我想现在连他们自己也承认这些是善，并且只能来自天主。但藉着\OriginalTerm{原质}{Hyle}，我指的是一种绝对无形、没有性质的质料，我们所感知的那些性质由之形成，正如古人所说。因为因此，木材在希腊文中也称为\OriginalTerm{木材}{\GreekText{ὕλη}}，因为它适用于工匠，并非它能制造什么，而是作为可供制作之物的材料。因此，那种\OriginalTerm{原质}{Hyle}也不应称为恶，它不能通过任何形象被感知，而几乎只能通过某种缺乏形象来思考。因为这种原质也具有接受形式的能力；如果它不能接受工匠所施加的形式，那么肯定也不可称之为质料。因此，如果形式是某种善（那些在形式方面卓越的人因而被称为美的，如同因容貌而被称为俊美），那么，形式的能力无疑也是某种善。正如智慧是善，无人怀疑能接受智慧也是善。而既然每种善都来自天主，那么无人应当怀疑，即使有质料，也只来自天主。\par

% structure: p0039 source=h2 target=section reason=relative-heading-rank; base=h2

\section{第十九章——真正存在是天主的专属特权}

因此，我们的天主庄严而神圣地对祂的仆人说：\BibleQuote{我是自有者}，以及\BibleQuote{你要对以色列子民说：那\Emph{自有者}打发我到你们这里来}（\BibleRef{出 3:14}）。因为祂真正地存在，因为祂是不变的。因为每一种变化都使不存在者成为存在；因此，那不变者才真正存在；但祂所造的其他一切事物，都从祂领受了存在，各自按自己的度量。因此，对于至高者，除了不存在之外，没有什么能与之相对；因此，正如一切善由祂而来，一切按本性存在的事物也由祂而来；因为一切按本性存在的事物都是善的。所以，每个本性都是善的，而一切善来自天主；因此，每个本性都来自天主。\par

% structure: p0041 source=h2 target=section reason=relative-heading-rank; base=h2

\section{第二十章——痛苦只存在于善的本性中}

但痛苦（有些人特别视之为恶），无论是在精神上还是在肉体上，都只能存在于善的本性中。因为任何存在物在痛苦中的抵抗本身，都涉及拒绝不再是它原本所是，因为它原本是某种善；但当一存在物被迫走向更好的状态时，痛苦是有益的；当被迫走向更糟的状态时，痛苦是无益的。因此，就精神而言，抗拒更大力量的意志引起痛苦；就肉体而言，抗拒更强壮身体的知觉引起痛苦。但无痛苦的恶更为恶劣：因为以不义为乐比哀叹败坏更糟；然而，甚至这种快乐也只能从获得低劣的善而产生。但不义是离弃更善的事物。同样，在身体中，有痛苦的伤口比无痛苦的腐烂更好，这种腐烂特别被称为腐败，主的死尸未曾见过，即未曾遭受，正如预言中所预言的：\BibleQuote{你必不让你的圣者见到腐朽}。因为谁能否认祂被钉子的刺伤所创伤，并被长矛所刺透？但即使人们特别称为身体腐败的东西，即腐烂本身，如果仍有某种东西可被消耗，它便藉着善的减损而加剧。但如果腐烂将彻底耗尽它，以至于没有善，则无本性存留，因为将没有任何东西可供腐烂来腐烂；因此，甚至腐烂也不再存在，因为根本没有它存在之处。\par

% structure: p0043 source=h2 target=section reason=relative-heading-rank; base=h2

\section{第二十一章——由尺度而言，事物被称为适度的大小。}

因此，如今按通常用法，细小卑微的事物也被说成具有尺度，因为其中仍保留着某种尺度，没有它，它们就不再是适度的大小，而是根本不存在。但那些因过度增长而被称为无度的事物，其受谴责正是由于过分；然而，这些事物本身也必须在天主之下受到某种约束，祂以度量、数目和重量安排了万有。\BibleRef{智 11:21}\par

% structure: p0045 source=h2 target=section reason=relative-heading-rank; base=h2

\section{第二十二章——尺度在某种意义上适用于天主本身}

但天主不能说有尺度，以免祂显得是有限的。然而，将尺度赐予万有、使它们得以存在的祂，并非无度者。同样，也不能称天主为有尺度，好像祂从谁那里领受了尺度。但若我们说祂是至高尺度，或许说得对；如果我们所谓的至高尺度是指至高美善的话。因为凡尺度，就其为尺度而言，都是善的；因此，没有尺度的事物不能被称作有尺度、谦抑、适可而止而不受称赞，尽管在另一意义上我们用尺度来表示界限，并称没有界限为无尺度，这有时用于赞美，如经上说：\Quote{祂的王国将没有界限。}\BibleRef{路 1:33} 因为也可说\Quote{没有尺度}，以尺度表示界限；因为统治没有尺度的，确实根本不统治。\par

% structure: p0047 source=h2 target=section reason=relative-heading-rank; base=h2

\section{第二十三章——何以有时可谈及恶的尺度、恶的形式、恶的秩序}

因此，恶的尺度、恶的形式、恶的秩序，之所以如此称呼，是因为它们小于应有，或是因为它们不适应于所当适应的事物；于是它们被称作恶，因为是相异和不协调的。例如，若有人说他没有以好的尺度行事，是因为他做得太少，或是在不该做的事上做了，或做得过多，或不合适；因此，那以恶的尺度所行之事之所以受责难，其唯一原因就是没有保持尺度。同样，形式被称为恶，或是与更俊美更美丽的事物相比，此形式较小，不是大小上，而是秀美上；或是因为它与其所应用的对象不和谐，显得相异和不合适。例如，若一个人裸露地走进公共场所，这种裸露若在浴室中看到则不构成冒犯。同样，秩序也被称为恶，当秩序本身在较低程度上保持时。因此，恶的不是秩序，而是失序；因为要么秩序过于不足，要么不如应有。然而，凡有尺度、形式、秩序之处，就有某种善和某种本性；但无尺度、无形式、无秩序之处，就无善、无本性。\par

% structure: p0049 source=h2 target=section reason=relative-heading-rank; base=h2

\section{第二十四章——由圣经的见证证明天主是不变的。天主子是受生而非受造的}

我们信仰所持守、理性以某种方式探究过的事，由神圣圣经的见证加以巩固，以便那些因悟性较弱而不能理解这些事的人，能相信神圣权威，从而得以认识。但那些明白却较少受教会文献训练的人，不要以为我们凭自己的理智提出这些事，而非来自那些书卷中所记载的。因此，关于天主是不变的，在圣咏中写道：\Quote{你要改变它们，它们便被改变；但祢本身是同一的。}并在智慧篇中论智慧说：\Quote{她虽自身不变，却使万象更新。}\BibleRef{智 7:27} 为此，保禄宗徒也说：\Quote{对那不可见的、不朽坏的、独一的天主。}\BibleRef{弟前 1:17} 雅各伯宗徒也说：\Quote{一切美好的赠与和完美的恩赐，都是从上而来，从光明之父降下，在祂那里没有变化或转动的阴影。}\BibleRef{雅 1:17} 同样，因为祂由自己生的即是祂自己，圣子自己曾简单地说：\Quote{我与父原是一体。}\BibleRef{若 10:30} 但圣子不是受造的，因为万有是藉着祂受造，经上如此记载：\Quote{在起初已有圣言，圣言与天主同在，圣言就是天主；这圣言在起初就与天主同在。万有是藉着祂造成的；凡受造的，没有一样不是藉着祂造成的。}\BibleRef{若 1:1}\par

% structure: p0051 source=h2 target=section reason=relative-heading-rank; base=h2

\section{第二十五章——最后这句话被一些人误解}

人狂言以为\Emph{无}应被理解为\Emph{某物}，且因\Quote{无}置于句末，便妄图迫人接受此等虚妄，对此无需理会。他们说：因此它被造成，且因被造成，无本身就是某物。他们出于反驳的热忱而失去理智，不明白说\Quote{没有祂，什么也没有被造成}或\Quote{没有祂，无被造成}并无区别。即使语序是后者，他们仍会说无本身就是某物，因为它是被造成的。因为对于真正是某物之事，若说\Quote{没有祂，一所房屋被造成}，只要理解没有祂造成了某物，即房屋，又有何区别？同样，因经上说\Quote{没有祂，无被造成}，既然无确实不是什么，则当真实而恰当地说出时，说\Quote{没有祂，无被造成}或\Quote{没有祂，什么也没有被造成}或\Quote{无被造成}并无区别。但谁愿与那些竟能对我这句话\Quote{它没有区别}说\Quote{因此它有些区别，因为无本身就是某物}的人交谈？然而脑筋清醒者视此为最明显的事：当说\Quote{它没有区别}时，应理解为某物，就如我说\Quote{它在任何方面都无关紧要}。但这些人若对某人说\Quote{你做了什么？}而他回答他什么也没做，则根据这种辩论方式，他们会错怪他说\Quote{因此你做了某事，因为你做了无；因为无本身就是某物}。他们也有主亲自将这字置于句末，当祂说：\Quote{我在暗地里什么也没有说过。}\BibleRef{若 18:20} 因此让他们阅读并保持沉默。\par

% structure: p0053 source=h2 target=section reason=relative-heading-rank; base=h2

\section{受造物由无而生}

因此，因为天主创造了祂并未从自身生出的一切，不是从已经存在的东西，而是从根本不存在的东西，即从无中造出，宗徒保禄说：\Quote{祂叫那不存在的，成为存在的。}\BibleRef{罗 4:17} 但更明确地记载于《玛加伯书》：\Quote{我恳求你，孩子，仰观天、地及其中万物，看看并知道，主天主造我们并非用这些东西。}\BibleRef{加下 7:28} 由此及诗篇所载：\Quote{祂发言，万物便造成。}显而易见，祂并未生出这些东西，而是以言语和命令创造了它们。凡非出自祂自己的，无疑出自无。因为没有任何先存之物可供祂创造，关于此点宗徒最明确地说：\Quote{因为万物出自祂、借着祂、并在祂内。}\BibleRef{罗 11:36}\par

% structure: p0055 source=h2 target=section reason=relative-heading-rank; base=h2

\section{\Quote{出自祂}与\Quote{来自祂}并非同一意思}

但\Quote{出自祂}与\Quote{来自祂}含义不同。凡来自祂的，可说出自祂；但并非一切出自祂的都正确地说来自祂。因为天地出自祂，因为祂造了它们；但不是来自祂，因为它们并非出自祂的实体。如同一个人生子并造房，儿子出自他，房屋出自他，但儿子来自他，房屋来自土和木。如此是因为作为人，他甚至不能从无中造出某物；但天主，万物出自祂、借祂、在祂内，无需任何非祂所造的材料来协助祂的全能。\par

% structure: p0057 source=h2 target=section reason=relative-heading-rank; base=h2

\section{罪并非来自天主，而是来自犯罪者的意志}

当我们听到\Quote{万物出自祂、借祂、在祂内}时，我们无疑应理解为所有自然存在的本性。因为罪，并不保存本性反而败坏本性，不是来自祂；圣经多次证明，罪来自犯罪者的意志，尤其在宗徒说：\Quote{啊！人哪，你判断做这样事的人，自己却做同样的事，你以为你能逃脱天主的审判吗？或是你轻视祂丰厚的仁慈、容忍和耐心，不知道天主的耐心是引导你悔改？但因你固执而不肯悔改的心，你在为自己积蓄忿怒，到忿怒的日子，即天主公义审判显现的日子，祂要照各人的行为报应各人。}\BibleRef{罗 2:3}\par

% structure: p0059 source=h2 target=section reason=relative-heading-rank; base=h2

\section{天主并不因我们的罪而受玷污}

尽管祂所建立的一切都在祂内，犯罪的人并不玷污祂；关于祂的智慧说：\Quote{她因纯洁无瑕，无所不通，没有任何污秽能浸入她。}\BibleRef{智 7:24} 因为我们必须相信，正如天主是不朽坏、不变的，同样也是不可玷污的。\par

% structure: p0061 source=h2 target=section reason=relative-heading-rank; base=h2

\section{美好之物，即使是最微小的、属地的，也都来自天主}

但天主也创造了最微小的事物，即属世和必朽之物，这无疑应从宗徒的那段话中理解，他论及我们肉身的肢体时说：\Quote{若一个肢体受光荣，所有肢体都一同欢乐；若一个肢体受苦，所有肢体都一同受苦。}接着又说：\Quote{天主将肢体一一安置在身体上，如祂所愿。}以及\Quote{天主调合了身体，对那缺乏的给予更多的荣耀，免得身体发生分裂，反而使肢体彼此同样照顾。}宗徒如此赞美肉身的度量、形式和秩序，你在所有动物，无论最大还是最小的肉身中都能找到；因为所有肉身都是属地之物中的，因此被视为最小的。\par

% structure: p0063 source=h2 target=section reason=relative-heading-rank; base=h2

\section{惩罚与赦罪均属于天主}

同样地，由于判定何种惩罚、多大惩罚适用每一个过犯，是属于天主的判断，而非人的判断，经上这样记载：\Quote{天主的智慧和知识何其深奥！祂的决断何其难测，祂的行径何其难寻！}\BibleRef{罗 11:33}同样地，由于藉着天主的恩慈，罪过得以赦免给悔改的人，\Person{基督}被派遣一事足以表明，祂并非以祂自己作为天主的本性，而是以我们的人性——这人性祂从一位女子所取得——为我们而死；这位宗徒这样陈明天主对我们所怀的恩慈与爱：\Quote{但天主对我们的爱在此显明了：当我们还是罪人的时候，\Person{基督}就为我们死了。何况现在我们已藉祂的血成义，更将藉祂从义怒中得救。因为如果我们尚与天主为敌时，已藉祂圣子的死而与祂和好，那么，在和好之后，更将藉祂的生命而得救。}\BibleRef{罗 5:8}但是，即便当应得的惩罚施加于罪人时，天主方面并无不义，他这样说：\Quote{我们可说什么呢？难道天主是那不义而施行义怒的吗？}但在另一处，他简短地告诫，恩慈与严厉皆出自祂，说：\Quote{你看，天主的恩慈与严厉：对跌倒的人，是严厉；对你，则是恩慈，只要你存留于恩慈中。}\par

% structure: p0065 source=h2 target=section reason=relative-heading-rank; base=h2

\section{第32章——害人的权能也出自天主}

同样地，甚至那些害人的权能也唯独出自天主，经上这样记载，智慧说：\Quote{藉着我，君王执政，暴君也藉着我占领土地。}\BibleRef{箴 8:15}宗徒也说：\Quote{没有权柄不是出于天主。}\BibleRef{罗 13:1}但是，这样的治理是应得的，这记载在约伯传中：\Quote{祂使伪君子作王，是因为百姓的顽梗。}关于以色列百姓，天主说：\Quote{我在怒气中将君王赐给他们。}\BibleRef{欧 13:11}因为并非不义：恶人获得害人的权能，既可以使善人的忍耐受到考验，又可以使恶人的不义受到惩罚。因为藉着给予魔鬼的权能，\Person{约伯}被考验以显明他的正义，\Person{伯多禄}被试探以免他自以为是，\Person{保禄}被击打以免他自高自大，\BibleRef{格后 12:7}\Person{犹达斯}被定罪以致自缢。\BibleRef{玛 27:5}因此，当天主藉着祂给与魔鬼的权能，公正地行了这一切事时，最终惩罚仍要归于魔鬼，不是因为他所行这些公正之事，而是因为他那属于他自己的、不义的害人意愿；那时将向那些执意与他的恶行同流合污的不敬者说：\Quote{你们这些被咒诅的，往那为魔鬼及其天使所预备的永火里去吧。}\BibleRef{玛 25:41}\par

% structure: p0067 source=h2 target=section reason=relative-heading-rank; base=h2

\section{第33章——恶天使并非由天主造成，而是由犯罪变恶}

但是，因为恶天使也并非由天主造成成为恶的，而是因犯罪而变恶，\Person{伯多禄}在他的书信中说：\Quote{天主既然没有宽容犯罪的天使，反把他们投入幽冥的地狱，交给他们黑牢，等候审判的惩罚。}\BibleRef{伯后 2:4}由此\Person{伯多禄}表明，最后的审判的刑罚仍等待着他们，关于这审判，主说：\Quote{你们这些被咒诅的，往那为魔鬼及其天使所预备的永火里去吧。}\BibleRef{玛 25:41}虽然他们已经刑罚性地接受了一个地狱，即一个充满烟气的下层空气作为监牢，但这空气既然也被称为天，并非那有星辰的天，而是这个较低的天，它的烟雾聚集云彩，飞鸟在此飞过；因为圣经既提到云雾的天，也称飞鸟为天上的。正如宗徒\Person{保禄}称呼那些恶天使——我们以虔敬的生活与之作战的仇敌——为\Quote{在天上界中邪恶的属灵事物。}\BibleRef{弗 6:12}为避免人理解这是指上层诸天，他在另一处清楚地说：\Quote{按照这世界的神——就是现今在悖逆之子身上运行的那位。}\par

% structure: p0069 source=h2 target=section reason=relative-heading-rank; base=h2

\section{第34章——罪不是追求邪恶的本性，而是离弃更好的本性}

同样地，因为罪或不义不是追求邪恶的本性，而是离弃更好的本性，这在圣经中这样记载：\Quote{天主所造的一切都美好。}\BibleRef{弟前 4:4}因此，天主在乐园中所栽种的每一棵树也确实是美好的。人触碰那禁止的树时，并不是追求邪恶的本性；而是因离弃更好的，犯下了恶行。既然造物主比祂所造的任何受造物都更好，就不应离弃祂的命令，而去触碰那被禁止的东西，无论它多么美好；因为离弃了更好的，反而追求受造物的美好，那受造物被触碰违背了造物主的命令。天主并没有在乐园中栽种邪恶的树；禁止人触碰它的是祂自己，祂是更好的。\par

% structure: p0071 source=h2 target=section reason=relative-heading-rank; base=h2

\section{第35章——那树被禁止于亚当，不是因为它邪恶，而是因为人服从天主是好的}

此外，祂作出这禁令，是为了表明：理性灵魂的本性不应囿于自身权能，而应服从于天主，并藉着服从保持其救恩的秩序，藉着不服从败坏它。因此，祂也称那棵祂禁止触碰的树为\Quote{知善恶树}\BibleRef{创 2:9}，因为当人违背禁令去触碰它时，他将经历罪的惩罚，从而知道服从的善与不服从的恶之间的区别。\par

% structure: p0073 source=h2 target=section reason=relative-heading-rank; base=h2

\section{第36章——天主的任何受造物都不是邪恶的，但滥用天主的受造物是邪恶的}

因为谁如此愚蠢，以致认为天主的受造物，尤其是种植在乐园中的受造物，是可责备的呢？事实上，连荆棘和蒺藜，就是大地依照天主公义的审判所产生，用以在劳作中折磨罪人的，也不应受到责备。因为即使这样的草本植物也有其尺度、形态和秩序，凡慎重考虑的人都会发现它们是值得赞美的；但它们对于那本性而言是恶的，这本性本应如此被抑制，作为罪的报应。因此，如我所说，罪不是追求一个恶的本性，而是离弃一个更好的本性，所以行为本身是恶的，而非那被罪人不当使用的本性是恶的。因为不当使用善的东西是恶的。因此，宗徒谴责某些人，称他们为神圣审判所定罪的人，\BibleQuote{他们敬拜事奉了受造物，而非造物主。}\BibleRef{罗 1:25} 他不谴责受造物（若有人这样做，便是对造物主不敬），而是谴责那些离弃更好的、不当使用善物的人。\par

% structure: p0075 source=h2 target=section reason=relative-heading-rank; base=h2

\section{第三十七章 天主善用罪人的恶行}

因此，如果所有本性都守护它们自己适当的尺度、形态和秩序，就不会有恶；但如果有人想要滥用这些善物，即使如此，他也不能战胜天主的旨意，天主知道如何公义地安排不义之事；以致如果他们自己通过其意志的不义而滥用祂的善物，祂便通过其德能的公义使用他们的恶行，正确地安排那些悖逆地将自己定意为罪的人受惩罚。\par

% structure: p0077 source=h2 target=section reason=relative-heading-rank; base=h2

\section{第三十八章 折磨恶人的永火并非邪恶}

因为永火本身，即用来折磨不敬神者的，并不是一个恶的本性，因为它有它的尺度、形态和秩序，没有被任何不义败坏；但它对于被惩罚者是一种恶的折磨，是由于他们的罪。同样，那边的光，因为它折磨患眼疾的人，也不是一个恶的本性。\par

% structure: p0079 source=h2 target=section reason=relative-heading-rank; base=h2

\section{第三十九章 火被称为永火，并非如天主那样，而是因为没有终结}

但是火是永久的，并非如天主是永久的，因为它虽无终结，却非无开始；但天主也是无开始的。其次，尽管它可被永久地用于惩罚罪人，它仍是可变化的本性。然而，那真正的永恒是真正的永存不朽，即那最高的不变性，完全不能被改变。因为，能够免于变化但变化可能是一回事，而完全不能变化是另一回事。因此，正如人被称为善，但不如天主那样善，关于天主曾说：\BibleQuote{除天主一个外，没有善的。}\BibleRef{谷 10:18} 也正如灵魂被称为不朽，但不如天主那样不朽，关于天主曾说：\BibleQuote{惟有祂独得不朽。}\BibleRef{弟前 6:16} 也正如人被称为智慧，但不如天主那样智慧，关于天主曾说：\BibleQuote{给那唯一智慧的天主。}\BibleRef{罗 16:27} 同样，火被称为永火，但不如天主那样，因为唯有祂本身是永存不朽和真正的永恒。\par

% structure: p0081 source=h2 target=section reason=relative-heading-rank; base=h2

\section{第四十章 天主不能受伤害，也无人能受伤害，除非出于天主公义的安排}

既然这些事如此，按照天主教信仰、健全的教义和对善解之人显明的真理，没有人能伤害天主的本性，天主的本性也不能不义地伤害任何人，或允许任何人伤害而免于受罚。\BibleQuote{因为作恶的人必受报应}，宗徒说，\BibleQuote{按他所做的恶；且天主不偏待人。}\BibleRef{哥 3:25}\par

% structure: p0083 source=h2 target=section reason=relative-heading-rank; base=h2

\section{第四十一章 摩尼教徒在恶的本性中放置了何等伟大的善，以及在善的本性中放置了何等伟大的恶}

但是，如果摩尼派愿意放弃对维护自己错误的这种有害热忱，并怀着对天主的敬畏而思考，他们就不会以最邪恶的方式亵渎，假设有两个本性：一个是善的，他们称之为天主；另一个是恶的，是天主所没有创造的。他们是如此错误，如此妄诞，甚至如此疯狂，以致他们看不见，即便在他们所谓的至恶的本性中，也安置了如此众多的善的事物：生命、能力、安全、记忆、理智、节制、德性、丰富、感觉、光明、温和、广延、数目、和平、尺度、形式、秩序；而在他们所谓的至善中，却安置了如此众多的恶的事物：死亡、疾病、遗忘、愚昧、混乱、无能、匮乏、迟钝、失明、痛苦、不义、耻辱、战争、放纵、畸形、乖谬。因为他们说，黑暗的众君王在自己的本性中也是活的，在自己的国度中是安全的，并且记忆和理解。因为他们说，黑暗之君以这样的方式发表演说，以至于没有记忆和理解，他既不可能说出这些话，也不可能被那些据说听到他说话的人所听见；而且他拥有适合其心灵和身体的性情，并凭着能力统治，关于他的元素有丰盛和多产，并且据说他们彼此感知，并感知到近处的光明，并且有眼睛可以远远看见光明；这些眼睛当然没有一些光明就不能看见光明（因此它们也被恰当地称为光明）；并且据说他们极其享受自己快乐的甜美，并且被有尺度的肢体和居所所限定。但除非那里有某种美，他们就不会爱自己的妻子，他们的身体也不会因肢体的配合而稳固；没有这些，那些摩尼派疯狂地说在那里发生过的事情就无法发生。除非那里有某种和平，他们就不会服从他们的君王。除非那里有尺度，他们就不会做任何事，除了饮食、狂怒或任何他们可能做的事，而没有丝毫的社会性；尽管连做这些事的人也不会具有确定的形状，除非尺度存在。但是现在摩尼派说他们做了那样的事，以至于不能否认他们在所有行动中都拥有适合自己的尺度。但如果那里没有形式，就没有自然性质会存留。但如果那里没有秩序，就不会有一些统治，另一些被统治；他们不会在自己的元素中和睦共处；最后，他们不会拥有适应于其位置的肢体，以至于他们不能做所有那些事，即摩尼派徒然虚构的事。但如果他们说天主的本性不会死，那么按照他们的虚妄，基督从死者中复活的是什么呢？如果他们说它不会生病，那么祂治愈的是什么呢？如果他们说它不会遗忘，那么祂提醒的是什么呢？如果他们说它不缺乏智慧，那么祂教导的是什么呢？如果他们说它不混乱，那么祂恢复的是什么呢？如果他们说它没有被征服和俘虏，那么祂解放的是什么呢？如果他们说它不匮乏，那么祂帮助的是什么呢？如果他们说它没有失去感觉，那么祂赋予生机的是什么呢？如果他们说它没有被弄瞎，那么祂照亮的是什么呢？如果它不痛苦，那么祂给予安慰的是什么呢？如果它并非不义，那么祂通过诫命纠正的是什么呢？如果它并非蒙羞，那么祂洁净的是什么呢？如果它并非处于战争，那么祂应许和平的是什么呢？如果它不缺乏节制，那么祂施加法律尺度的是什么呢？如果它并非畸形，那么祂改造的是什么呢？如果它并非乖谬，那么祂修正的是什么呢？因为所有这些由基督做成的事，他们说，都不是归于那由天主所造、并因自己的自由意志在犯罪中败坏了的事物，而是归于天主的本性本身，甚至归于天主的实体本身，即天主自己所是。\par

% structure: p0085 source=h2 target=section reason=relative-heading-rank; base=h2

\section{第42章——摩尼派关于天主本性的亵渎之词}

什么能与这些亵渎相比？绝对没有，除非考虑其他派别的错误；但如果将那个错误与它自身尚未说到的另一方面相比，它将被定为更恶劣、更可憎的亵渎。因为有些人说，某些灵魂（他们坚持认为这些灵魂与天主的本体相同，且完全同质）并非自愿犯罪，而是被黑暗种族（他们称之为邪恶）战胜和压制；为了与这黑暗种族作战，他们并非自愿降下，而是奉圣父命降下，却被永远锁在可怖的黑暗之球中。因此，按照他们亵渎神明的胡言，天主在某一部分将自己从一个巨大邪恶中解放出来，但在另一部分（他无法解放的部分）却再次定自己的罪，并且好像从上方战胜了仇敌。哦，可耻的、难以置信的胆量，竟敢如此相信、谈论、宣告关于天主的事！当他们试图辩护时（为闭着眼睛盲目冲入更糟的境地），他们说邪恶本性的混杂造成了这些，以便天主的善性遭受如此大的恶：因为这善性在其自身领域内不可能也不会遭受其中任何一点。仿佛一种本性因其不伤害自己而被称赞为不朽坏，而不是因为它不能受伤害于其他。那么，如果天主的本性伤害黑暗的本性，而黑暗的本性也伤害天主的本性，那么就有两件邪恶的事物在相互伤害，而黑暗种族的态度更好，因为它即使造成伤害也是不情愿的；因为它不想造成伤害，而是想享受属于天主的善。但天主却想要灭绝它，正如摩尼在他那封关于毁灭的\WorkTitle{基础}的书信中最公开地疯狂咆哮。因为他忘了不久之前曾说过：\Quote{但他最灿烂的国度如此建立在光明幸福之地上，以至于永不能被任何人动摇或震动；}随后他又说：\Quote{最蒙福光明之圣父知道那将从黑暗中兴起的巨大毁灭和荒凉将威胁他的神圣世界，除非他送出一位卓越、显赫、大有能力的神明与之对抗，藉此他同时战胜并摧毁黑暗种族；待其被灭绝后，光明居民将享受永恒的安息。}看，他惧怕威胁他世界的毁灭和荒凉！那些国度确实建立在光明幸福之地上，以至于永不能被任何人动摇或震动？看，出于恐惧，他想要伤害邻近的种族，他努力要毁灭和灭绝它，以便光明居民享受永恒的安息。他为何不加上\Quote{永久的捆绑}呢？这些被他永远锁在黑暗之球中的灵魂，不就是光明居民吗？他明确说到他们\Quote{容许自己偏离先前光明的本性？}当他被迫承认他们因自由意志犯罪时，却违背他的意愿，想把罪仅仅归因于相反本性的必然；他处处不知该说什么，仿佛他自己已在他捏造的黑暗之球中，寻找却找不到逃脱之路。但让他随意对受迷惑的可怜人说什么吧，这些人尊崇他远超过基督，他以此代价向他们兜售如此漫长而亵渎的寓言。让他随意说什么吧，让他把黑暗种族仿佛锁在一个球体里（如同在监狱中），并把光明的本性系在外面（他许诺在不死之敌灭绝后给予永恒的安息）：看，光明的惩罚比黑暗的更重；神性本体的惩罚比敌对种族更重。因为虽然后者在黑暗之中，但它的本性就是居住在黑暗中；而灵魂与天主完全相同，他说它们不能被接纳进入那些平安的领域，被隔绝于圣光的生命和自由，被锁在之前提到的可怖球体中。由此他说：\Quote{这些灵魂将被遗留在同一黑暗之球中，执着于他们所爱的事物，因自己的功过而招致如此。}这难道不是确定的自由意志选择吗？看他多么疯狂地忽视自己所说的话，用自相矛盾的陈述展开一场比对抗黑暗种族之神更糟糕的战争。因此，如果光明的灵魂因爱黑暗而被定罪，那么爱光明的黑暗种族就被不公正地定罪了。而且黑暗种族从一开始就爱光明，或许暴力地爱，但仍是希望占有它，而非灭绝它；而光明的本性却想在战争中灭绝黑暗；因此当它失败时，它爱了黑暗。你选择吧：它是被必然所迫爱黑暗，还是被自由意志诱惑？如果出于必然，为何被定罪？如果出于自由意志，为何天主的本性陷入如此大的不义？如果天主的本性被必然所迫爱黑暗，那么它没有战胜，而是被战胜了；如果出于自由意志，为何这些可怜虫还要犹豫，不敢将犯罪意志归因于天主从虚无中造出的本性，以免他们将此归因于他所生的光？\par

% structure: p0087 source=h2 target=section reason=relative-heading-rank; base=h2

\section{第四十三章——在与邪恶混杂之前，许多邪恶被摩尼派归诸天主的本性}

我们若还要指出，在他们最疯狂地相信的那个愚昧荒诞故事中，在罪恶混杂之前，他们所谓的\OriginalTerm{光明本性}{nature of light}中就已经存在巨大的罪恶，那又当如何呢？在这些亵渎之上还能再加添什么吗？因为在冲突之前，就已经有严酷且不可避免的战斗必然性：这诚然是一个巨大的罪恶，发生在罪恶与善混杂之前。让他们说说，这从何而来，当时尚没有混杂发生？但若没有必然性，那就有自由意志：从何而来这如此巨大的罪恶，竟使天主自己愿意伤害自己的本性（这本性本是敌人无法伤害的），把它送去残酷地混杂、卑劣地净化、不义地受罚？看啊，在敌对本性的任何罪恶与之混杂之前，就已经有了恶毒、有害、野蛮意志的巨大罪恶！或者，或许他并不知道这事会临到他的肢体身上，即\Quote{它们应该爱慕黑暗，对圣洁的光明怀有敌意，正如\Person{摩尼}所说，即不仅对它们自己的天主，而且也对那赋予它们存在的圣父}？因此，在黑暗本性的任何罪恶与之混杂之前，这无知的巨大罪恶从何而来？但若他知道这事会发生，那么，若他不为将要发生的自己本性的污染和受罚而忧伤，他就有永久的残忍；若他忧伤，他就有永久的痛苦：从何而来你们至尊之善在与其至尊之恶混杂之前的这巨大罪恶？确实，那本性中自身被束缚在那领域的永恒锁链中的部分，若它不知道这一命运等待它，那么在天主本性中就有永久的无知；但若它知道，则有永久的痛苦：在敌对本性的任何罪恶与之混杂之前，这巨大罪恶从何而来？或者，或许它在其伟大的\OriginalTerm{爱德}{charity}中，因着它的惩罚，为光明居民中余下的人预备了永恒的安息而欢喜？谁若看出说这话是何等可憎，就让他发出弃绝。但若这样做是为了至少善的本性本身不致对光明怀有敌意，那么有可能，或许不是天主的本性，而是仿佛某个人，可被视为可称赞的，他为了自己的国家愿意承受某种恶，而这恶只能是一时的，而非永远的；但现在他们还说，那在黑暗领域中的束缚是永恒的，并且不是某样东西，而是天主的本性；如果天主的本性欢喜它应爱慕黑暗，对圣洁的光明怀有敌意，那么这确实是一种最不义、最可憎、无法言说地亵渎的喜乐。在敌对本性的任何罪恶与之混杂之前，这如此怪异可憎的罪恶从何而来？谁能忍受如此反常和不虔敬的疯狂，竟将如此巨大的善归于至恶，又将如此巨大的恶归于至善——即天主？\par

% structure: p0089 source=h2 target=section reason=relative-heading-rank; base=h2

\section{第四十四章——\Person{摩尼}所想像的天主中难以置信的丑恶}

但是现在，当他们谈论天主本性的那一部分，说它普遍混合在天上、地上、一切干湿的身体中，各种血肉中，树木、草木、人类和动物的一切种子中时——并非像我们论及天主那样，以其神性的权能临在，为治理和统治万物，无玷污、无侵害、无腐化，不与它们有任何联系；而是被束缚、压迫、玷污，需要被解开和解放，据他们说，不仅通过太阳和月亮的来回运行，以及通过光明的力量，而且还通过他们的选民——这种错误所向他们推荐的，即使没有诱导他们接受，但那些亵渎且难以置信的丑恶行径，说出来都令人恐怖。因为他们说，光明的力量变成俊美的男性，被置于黑暗种族的妇女对面；又同样，这些力量变成美丽的女性，被置于黑暗种族的男性对面；藉着他们的美丽，他们点燃了黑暗王子们最污秽的情欲，而藉此，生命实质——即天主的本性，据他们说被束缚在这些王子的身体中——因着情欲而松弛的肢体中解脱出来，飞逸而去，当它被吸收或净化后，便得解放。这些可怜的人诵读这些，说这些，听这些，信这些，并如下所述，记载在他们的\WorkTitle{宝库}第七卷中（他们如此称呼\Person{摩尼}的某一著作，这些亵渎之言即写在其中）：\Quote{那时，蒙福的圣父，他拥有光明之船、小室、居所或大小不同之物，按照他内住的仁慈，带来帮助，使他的生命实质从邪恶的束缚、狭窄和折磨中被引出并解放。因此，藉着他不可见的点头，他转变那些力量（即被持守在这艘最辉煌之船中的他的力量），并使它们产生出敌对的力量，这些敌对的力量已被安排在天空的各个区域。既然这些力量由两性构成，即男性和女性，他命令上述力量一部分以无胡须青年的形象出现，对付女性的敌对种族，一部分以光明少女的形象出现，对付男性的敌对种族；因为他知道所有这些敌对力量，由于它们内在的致命且最污秽的情欲，非常容易被俘虏，向这些显现的最美丽形象投降，从而被分解。但是你可以知道，我们这位同样的蒙福圣父与他的力量是同一的，他为了必要的理由将这些力量转变为无玷污的青年和少女的相似形。但他使用这些力量作为自己的武器，并通过它们完成他的旨意。但是有光明之船满载这些神圣力量，它们被安置在类似婚姻的形象中，面对地狱种族，并迅速而轻松地实现它们所计划的事。因此，当理性要求这些同样的神圣力量向男性显现时，他们立刻也通过他们的服饰显示出最美丽少女的形像。然后再处理女性时，放下少女的形像，他们显示出无胡须青年的形像。但藉着他们这英俊的外表，热情和情欲增长，这样他们最坏思想的链条被解开，那被他们的肢体所持守的活灵，藉此机会松弛而逃脱，并与它自己最纯净的空气混合；当灵魂彻底净化后，上升到光明之船，这些船已被预备好以运送他们并摆渡他们到自己的家乡。但那些仍然带有敌对种族污点的灵魂，逐渐通过波涛和火焰下降，与树木和其他植物以及一切种子混合，并被投入各种火焰中。并且，从那艘伟大而最光荣的船中，青年和少女的形象如何向居住在天空中并具有火性本性的敌对力量显现；并且从那英俊的外表中，被持守在他们肢体中的部分生命被释放，通过火焰被引到地上：同样，那艘住在活水之船中的至高力量，以青年和圣洁少女的形像向那些本性寒冷潮湿、并排列在天空中的力量显现。实际上，对于其中的女性，显现出青年的形像；但对于男性，显现出少女的形像。藉着他神圣而最美丽位格的转变和多样性，潮湿寒冷种族的男女王子们被解开，他们内在的生命成分逃逸；但任何剩余的，一旦松弛，便通过寒冷被引到地上，并与一切黑暗种族混合。}谁能忍受这个？谁能相信——不是相信它是真的，而是相信甚至有人能说出这样的话？看哪，那些害怕谴责教导这些事的\Person{摩尼}的人，却不害怕相信一位做着这些事并受着这些苦的天主！\par

% structure: p0091 source=h2 target=section reason=relative-heading-rank; base=h2

\section{第四十五章．论某些关于摩尼教徒自身、并非无理由相信的不可言喻的可耻之事}

但他们说，通过他们自己的选民，那同一混合的部分和天主的本性被净化，即通过饮食（因为他们说它被束缚在一切食物中）；当这些被选民取用作为身体营养的饮食时，它通过他们的圣洁被解开、封印并解放。这些可怜的人没有注意到，关于他们有这样的信念并非没有充足理由，他们徒然否认，只要他们不谴责\Person{摩尼}的书籍并停止做摩尼教徒。因为如果，如他们所说，天主的一部分被束缚在一切种子中，并通过选民的饮食被净化；那么，谁不能合理地相信，他们所做的正是他们在\WorkTitle{宝库}中所读到的、发生在天上力量和黑暗王子之间的事？因为他们确实说他们的肉体也来自黑暗种族，并且他们毫不犹豫地相信并肯定，束缚在他们内的生命实质是天主的一部分。这如果确实要通过饮食被解开和净化，正如他们可悲的错误迫使他们承认的那样；那么，谁看不见，谁不为随之而来的巨大和不可言喻之事而战栗？\par

% structure: p0093 source=h2 target=section reason=relative-heading-rank; base=h2

\section{第四十六章——\WorkTitle{基本书信}的不可言说的教义}

因为他们甚至说，第一个人亚当是由某些黑暗的王子创造的，以便光能被他们控制，以免它逃脱。因为在那封他们称为\WorkTitle{基本书信}的书信中，\Person{摩尼}写到如下内容，关于黑暗的王子（他们将其描述为第一个人的父亲）如何对他的其他同盟黑暗王子说话，以及他如何行事：\Quote{因此，他用邪恶的发明对在场者说：\InnerQuote{你们认为这个正在升起的巨大光是什么？看，天极如何移动，它如何震动大部分权能。因此，我最好先向你们索要你们权能中的任何光；因为这样我将塑造那位在荣耀中出现的那位伟大者的形象，通过这个形象我们得以统治，在某种程度上摆脱黑暗的交谈。}闻听这些，并在彼此间长久商议后，他们认为提供所要求的东西是最公正的。因为他们没有信心永远拥有他们所有的光；因此他们认为最好将它献给他们的王子，绝没有希望这样他们就能统治。因此必须考虑他们如何提供他们所拥有的光。因为这也散布在所有神圣的圣经和天上的奥秘中；但对于智者来说，知道它如何被给予是相当容易的：因为真正并忠实地愿意考虑的人会立即公开地知道它。既然聚集在一起的是一群混杂的人群，有男有女，他促使他们彼此交合：在这种交合中，男性射出种子，女性怀孕。但后代像那些生养他们的人，第一个获得父母力量的似乎最大的部分。将这些作为特别的礼物，他们的王子欢欣。正如我们现在甚至看到的，邪恶的本性从那里夺取力量，形成身体的创造者，同样，上述王子，拿走了他同伴的后代（这些后代在出生过程中与父母同时具有他们的感觉、敏锐、光），吞吃了它们；从这种食物中获取了许多权能，其中不仅有刚毅，而且还有来自祖先凶残种族的更多的狡猾和堕落的情感，他召唤自己的配偶（与他同源）到他自己这里，像其他人一样，他将所吞吃的所有邪恶的丰富排出，他自己也添加了一些来自他自己的思想和权能，这样他的性情成为他所倾泻的一切事物的形成者和安排者；其配偶接受了这些，就像最优良的土壤习惯于接受种子。因为在其中，所有天上和地下权能的形象被构建和编织在一起，以便所形成的东西获得了所谓完整圆球的相似。}\par

% structure: p0095 source=h2 target=section reason=relative-heading-rank; base=h2

\section{第四十七章——他强迫人犯下可怕的丑行}

啊，可憎的怪物！啊，被欺骗灵魂的可咒的沉沦与毁灭！我不是在说关于天主本性的这种被如此束缚的言论的亵渎。让那些被致命错误欺骗和追逐的可怜人至少注意这一点：如果他们天主的一部分被男女交合所束缚（他们声称通过吃这来解放和净化），那么这种不可言说的错误的必要性迫使他们不仅要从面包、蔬菜和水果（他们公开参与的似乎只有这些）中解放和净化天主的部分，而且也要从可能通过交合而被束缚的东西中解放和净化，如果受孕发生的话。据说有些人不仅在\Place{帕夫拉戈尼亚}，而且在高卢，在公开的法庭上承认他们这样做，正如我在罗马从一位天主教基督徒那里听到的；当被问到他们根据什么著作的权威做这些事时，他们泄露了关于我刚才提到的\WorkTitle{宝库}的这件事。但当有人指责他们这一点时，他们习惯于回答说，某个敌人从他们的行列中，即从他们的选民行列中，退出并制造了分裂，创立了这种最污秽的异端。由此可见，即使他们自己不实行此事，一些实行此事的人也是根据他们的书籍来做的。因此，如果他们厌恶这罪行，就让他们抛弃这些书籍（如果他们坚持这些书籍，他们就被迫犯这罪）；或者，如果他们不犯，他们就努力违背书籍而活得更纯洁。但当有人对他们说：\Quote{要么从你能得到的任何种子中净化光，这样你就不能拒绝做你声称不做的事；要么就诅咒\Person{摩尼}，因他说天主的一部分在一切种子中，且它被交合所束缚，但任何光，即上述天主的部分，如果成为选民的粮食，就会通过被吃而得到净化。}这时他们做什么呢？你们看到他迫使你们相信什么，你们还犹豫要诅咒他吗？我是说，当有人这样对他们说时，他们做什么？他们逃到什么样的借口里，要么诅咒如此邪恶的教义，要么犯下如此邪恶的丑行？相比之下，我已经提请注意的所有那些不可容忍的邪恶似乎都可容忍了，即他们说天主的本性被需要所迫而发动战争，它要么因永久的无知而安稳，要么因永久的悲伤和恐惧而烦扰，当混杂的败坏和永罚的锁链临到它时，最终由于冲突而被俘虏、压抑、玷污，在虚假的胜利之后被永远束缚在一个可怕的球体中，与它原初的福乐分离；而这些如果就其本身考虑，是无法忍受的？\par

% structure: p0097 source=h2 target=section reason=relative-heading-rank; base=h2

\section{第四十八章——\Person{奥古斯丁}祈求摩尼教徒恢复理智}

主，祢的忍耐何其伟大，祢充满慈悲和恩惠，缓于发怒，富于仁慈和信实；\BibleQuote{祢使太阳上升，光照善人和恶人，又降雨给义人和不义的人；}\BibleRef{玛 5:45}\BibleQuote{祢不愿恶人死亡，却愿他回头并获得生命；}\BibleRef{则 33:11}\BibleQuote{祢逐步责罚，却给人留有悔改的余地，使他们摒弃邪恶，信靠祢，主；}\BibleRef{智 12:2}\BibleQuote{祢以忍耐引导人悔改，尽管许多人按着他们刚硬和不悔改的心，为自己积攒愤怒，直到愤怒和祢公义审判显现的日子，祢必按各人的行为报应各人；}\BibleRef{罗 2:4}\BibleQuote{当一个人转离自己的邪恶，归向祢的仁慈和真理的那一天，祢将忘记他的一切过犯：}\BibleRef{则 18:21}求祢站在我们面前，恩赐我们，藉着我们的职务——祢曾乐意藉此驳斥这可憎且极其骇人的谬误——使许多人已经得到释放，也让许多人得释放，无论是藉着祢的圣洗圣事，还是藉着痛悔和谦卑心灵的祭献，在忏悔的忧苦中，他们得以堪当领受因无知而冒犯祢的罪恶和亵渎的赦免。因为除了祢的卓绝仁慈和权能，以及祢圣洗的真理和祢神圣教会中天国的钥匙之外，一无所用；因此，我们决不该对人绝望，只要他们还在世上存活，在祢的忍耐之下；他们即使知道对祢存有或说出这样的想法是何等大恶，却因使用或获得现世或地上的利益而被囚禁在那邪恶的教门中，若他们受祢的责备，以某种方式逃向祢那不可言喻的良善，并选择天上永恒的生命，胜过一切肉欲生活的诱惑。\par

\SourceRecord{Translated by Albert H. Newman.；Nicene and Post-Nicene Fathers, First Series；Vol. 4.；Edited by Philip Schaff.；Buffalo, NY: Christian Literature Publishing Co.,；1887.；<http://www.newadvent.org/fathers/1407.htm>.}
